% Part I, chapter 5: Header metadata.
\chapter{Header metadata}
\label{chap:metadata}

Many datasets store facts about a recording in the info strings of its
header (\cref{sec:info-strings}), as \code{Key: value} lines, for example:
\begin{lstlisting}[style=text]
# Age: 63
# Sex: F
# Dx: 164889003
\end{lstlisting}
These entries are the \emph{metadata} of the record. Five functions of the
high-level API, added in version 1.2.0, read and change them without
changing anything else in the header: \code{read\_metadata},
\code{get\_metadata}, \code{set\_metadata}, \code{remove\_metadata} and
\code{update\_metadata}. \code{read\_header} also returns the entries, and
\code{write\_record} writes them. The C package has no counterpart. The
library functions \code{putinfo} and \code{setinfo}, and the tools, are
unchanged, and so is \code{getinfo}, except for the rule by which it finds
the record line (\cref{sec:info-strings}), which the metadata functions
share; the C library reads the headers that these functions change exactly
as before (\file{tests/metadata/check\_c.py} checks this with the original
programs). The seven implementations follow the
rules of this chapter identically: they pass the same cases
(\file{tests/metadata/}) and write byte-identical headers.

\section{Entries}
\label{sec:metadata-entries}

A header is split into lines after each line feed (LF). Each line has a line
ending: CR LF, LF, or none (the last line of a file that does not end with a
line feed). The \emph{content} of a line is the line without its ending.

\paragraph{Info strings.} The record line is found as the header parser of
the library (\code{readheader}) finds it: it is the first line whose
content, up to a NUL character if there is one, has a token (separated by
spaces, tabs, carriage returns and line feeds) that does not begin with
\code{\#}; in other words, the first line that holds something other than
white space and does not begin, after white space, with \code{\#}. The
lines before it (comments, and empty or blank lines) are not info strings,
and are never read or changed. Every line after the record line whose
content begins with \code{\#} is an info string, as \code{getinfo} reads a
header (so a comment between two signal lines is an info string too). A
header with no record line has no info strings, and changing its metadata is
an argument error. This rule is the same in
every mode, legacy mode included (since version 1.2.2; before, the record
line was the first line whose content does not begin with \code{\#}, the
rule of the C function \code{getinfo}, which takes an empty first line for
the record line).

\paragraph{Entries and free comments.} An info string is an \emph{entry} if
the text after its \code{\#} contains a colon. The \emph{key} is the text
before the first colon, and the \emph{value} is the text after it; spaces,
tabs, vertical tabs, form feeds and carriage returns are removed at both ends
of each. If the key is then empty, or contains \code{\#}, the info string is
a \emph{free comment}. Free comments, and info strings without a colon, are
never returned and never changed. So \code{\#Age:45}, \code{\# Age: 45} and
\code{\#~~Age~:~45} all give the key \code{Age} and the value \code{45}, and
\cref{tab:metadata-reading} shows a few other cases.

\begin{table}[htbp]
\caption{How info strings are read.}
\label{tab:metadata-reading}
\begin{tabularx}{\linewidth}{@{}P{5.6cm}L@{}}
\toprule
{\tabhead Info string} & {\tabhead Entry (key, value)} \\
\midrule
\code{\#Age: 63} & (\code{Age}, \code{63}) \\
\code{\# Sex:F} & (\code{Sex}, \code{F}) \\
\code{\#~~Dx : 164889003, 59118001~} & (\code{Dx},
  \code{164889003, 59118001}) \\
\code{\#Notes:} & (\code{Notes}, the empty string) \\
\code{\#Time: 12:30:00} & (\code{Time}, \code{12:30:00}): the key ends at the
  first colon \\
\code{\#\textless{}tab\textgreater{}Tab key\textless{}tab\textgreater{}:%
  \textless{}tab\textgreater{}value with tabs} & (\code{Tab key},
  \code{value with tabs}) \\
\code{\# A free comment without a colon} & none: a free comment \\
\code{\#: empty key} & none: the key is empty \\
\code{\#\# Hash: not an entry} & none: the key (\code{\# Hash}) contains
  \code{\#} \\
\code{\#Produced by xform from record 100, beginning at 0:0} &
  (\code{Produced by xform from record 100, beginning at 0}, \code{0}) \\
\bottomrule
\end{tabularx}
\end{table}

The last line of \cref{tab:metadata-reading} is an info string of the sample
record \code{100s}: free text that happens to contain a colon is read as an
entry. Such an entry is harmless: it is changed only if its key is named.

\paragraph{Keys} are compared without regard to letter case, for the letters
A--Z and a--z only (no other character is changed: \code{\'E} and
\code{\'e} differ). So \code{get\_metadata(record, "age")} finds the entry
\code{Age}. Entries are returned in file order.

\paragraph{Duplicate keys.} A header should not hold two entries whose keys
are equal, ignoring letter case. If one does (a file written by other
software), the readers (\code{read\_metadata}, \code{get\_metadata},
\code{read\_header}) return the first of them, and give one warning for each
such key, naming the line numbers of its entries (from 1), such as
\code{header of record dups has more than one entry for key age (lines 4,
6); the first is used}; \code{set\_metadata} and \code{update\_metadata}
refuse to change that key (an argument error), and \code{remove\_metadata}
removes all of them, which repairs the header.

\section{Arguments}
\label{sec:metadata-arguments}

Before a key or a value given to a function is checked, the spaces, tabs,
vertical tabs, form feeds and carriage returns at both ends are removed.
\begin{itemize}
\item A \textbf{key} must then not be empty, and must not contain \code{:},
  \code{\#}, a carriage return, a line feed or a NUL character. It may
  contain spaces, as in \code{Date of birth}.
\item A \textbf{value} is text, which may be a whole sentence. It must not
  contain a carriage return, a line feed or a NUL character; it may be
  empty.
\item Text is one byte per character (as in the rest of \wfdbx{}: in
  Python, Latin-1 strings, or bytes as given).
\item In \code{update\_metadata}, no two keys may be equal, ignoring letter
  case, and a value of none (\code{None}, \code{nothing}, \code{NULL},
  \code{null}, or a missing value, depending on the language) means
  ``remove''.
\item The record name \code{-} (the standard input) is an argument error,
  and so is \code{get\_metadata} with an invalid key.
\end{itemize}
Every argument is checked before the header is read, and every change
before anything is written: after an error, the header is unchanged.

\section{Writing}
\label{sec:metadata-writing}

\paragraph{The written form.} An entry is written as \code{\# Key: value}:
\code{\#}, a space, the key as the caller spelled it, \code{:}, and, if the
value is not empty, a space and the value. An entry with an empty value is
therefore written \code{\# Key:}, with no trailing space.

\paragraph{The changes.}
\begin{itemize}
\item \code{set\_metadata}: if the header has one entry for the key, its line
  is replaced, at the same position, keeping its line ending, even if the
  value is the same (the line is then rewritten in the written form, with the
  caller's spelling of the key). Otherwise a new line is inserted.
\item \code{remove\_metadata}: every line of an entry for the key is deleted,
  with its line ending. The function returns the number of lines deleted (0
  if there is none).
\item \code{update\_metadata}: the changes are applied in the order given,
  each as \code{set\_metadata} or \code{remove\_metadata} would apply it, and
  written at once.
\end{itemize}

\paragraph{Where new lines go.} A new line is inserted after the last info
string of the header, or, if the header has none, after its last line whose
content is not blank (not made only of spaces, tabs, vertical tabs, form
feeds and carriage returns). Its line ending is that of the record line, or,
if the record line has none, that of the first line that has one, or CR LF
(the line ending of the headers written by the WFDB library). If the line
before it has no line ending, that line gets the same ending.

\paragraph{What is kept.} Every other byte of the header is kept: the
comments and the empty or blank lines before the record line, the record
line, the signal lines and their checksums, the free comments, and the line
endings. If nothing changes (and \code{dir} is not given), nothing is
written.

\paragraph{Atomic replacement.} The new header is written to a temporary
file in the directory of the destination, which is never more open than the
destination (in Python, it is created readable and writable by its owner
only). The temporary file then gets the permissions of the destination (or,
for a new file in \code{dir}, those of a new file) and replaces it (by
renaming it). An error never leaves a partly written header, and no other
file is left in the directory.

\section{Which file}
\label{sec:metadata-file}

The header is opened as the library opens the header of a record, after
removing a trailing \file{.hea} from the name: under the input directory, or
as given for an absolute path, a \code{\textasciitilde} path or a URL (a URL
only if network access is on); see \cref{sec:record-paths}. The functions
take the options of the high-level readers for this, an input directory
(\code{input\_dir}) and a network permission (\code{network}), for the
duration of the call. A header that cannot be opened is a file error.

A change is written to the file that was read, unless the call gives an
output directory \code{dir}: the changed header is then written to the file
of the same name in that directory, which must exist (else a file error),
and the header that was read is unchanged. The signal files are not copied.
A header read from a URL cannot be changed without \code{dir} (an argument
error).

\paragraph{Multi-segment records.} Only the named header is read or
changed: for a multi-segment record, that is the master header, unless the
name of a segment is given. (The info strings of a master header are read by
\code{getinfo} like those of any other header.)

\paragraph{Info files.} Only the header file (\file{.hea}) is read, never
the info file (\file{.info}) whose comment lines \code{getinfo} also
returns.

\section{read\_header and write\_record}
\label{sec:metadata-header}

\code{read\_header} returns the entries of a header, as \code{read\_metadata}
returns them (from the \file{.hea} file only), as the \code{metadata} of its
result, beside its \code{comments}, which are all the info strings.

\code{write\_record} writes the entries given in its \code{metadata}
argument (a mapping from keys to values, or a list of key and value pairs,
depending on the language) after its \code{comments}, in the written form
\code{\# Key: value}, through \code{putinfo}. Its keys must be unique,
ignoring letter case, and none may be the key of an entry in
\code{comments}; a value of none is an argument error. This is checked
before any file is created.

\section{An example}
\label{sec:metadata-example}

The header \file{tests/metadata/headers/basic.hea} (with CR LF line endings)
is:
\begin{lstlisting}[style=text]
# Converted on: 2026-11-01 (a comment before the record line)
basic 2 360 21600
100s.dat 212 200 11 1024 995 21537 0 MLII
100s.dat 212 200 11 1024 1011 -3962 0 V5
#Age: 63
# Sex:F
#  Dx : 164889003, 59118001
# A free comment without a colon
#Notes:
# Recorded by: Holter monitor, model X
\end{lstlisting}
\code{read\_metadata("basic")} returns the five entries, in file order:
(\code{Age}, \code{63}), (\code{Sex}, \code{F}), (\code{Dx},
\code{164889003, 59118001}), (\code{Notes}, the empty string) and
(\code{Recorded by}, \code{Holter monitor, model X}). The first line is
before the record line, and is not an info string; the free comment is not
an entry. In Python, the changes
\begin{lstlisting}[language=Python]
wfdbx.update_metadata("basic", [("age", "64"), ("Sex", None),
                                ("Weight", "72"), ("Height", "170")])
\end{lstlisting}
replace the entry \code{Age} (at its position, with the caller's spelling
\code{age}), remove the entry \code{Sex}, and add two new entries after the
last info string, with the line ending of the record line (CR LF). The other
lines are unchanged, byte for byte, including the unusual spacing of the
\code{Dx} line. The header becomes
\file{tests/metadata/expected/update\_basic.hea}:
\begin{lstlisting}[style=text]
# Converted on: 2026-11-01 (a comment before the record line)
basic 2 360 21600
100s.dat 212 200 11 1024 995 21537 0 MLII
100s.dat 212 200 11 1024 1011 -3962 0 V5
# age: 64
#  Dx : 164889003, 59118001
# A free comment without a colon
#Notes:
# Recorded by: Holter monitor, model X
# Weight: 72
# Height: 170
\end{lstlisting}

A header whose first line is empty, such as
\file{tests/metadata/headers/leading.hea}:
\begin{lstlisting}[style=text]
(an empty line)
# Age: preamble
leading 2 360 21600
100s.dat 212 200 11 1024 995 21537 0 MLII
100s.dat 212 200 11 1024 1011 -3962 0 V5
# Age: actual
\end{lstlisting}
has one entry, (\code{Age}, \code{actual}): its record line is the third
line, as for the header parser, and \code{\# Age: preamble}, before it, is
not an info string. \code{set\_metadata("leading", "Sex", "F")} adds
\code{\# Sex: F} after \code{\# Age: actual}, and leaves the first two lines
as they are. In Python:
\begin{lstlisting}[language=Python,style=example]
import wfdbx
with open("leading.hea", "w") as f:
    f.write("\n# Age: preamble\nleading 2 360 21600\n"
            "100s.dat 212 200 11 1024 995 21537 0 MLII\n"
            "100s.dat 212 200 11 1024 1011 -3962 0 V5\n# Age: actual\n")
print(wfdbx.read_metadata("leading"))           # [('Age', 'actual')]
print(wfdbx.read_header("leading").comments)    # [' Age: actual']
\end{lstlisting}

The following Python example writes a record with metadata, from the first
second of record \code{100s}, and changes its metadata:
\begin{lstlisting}[language=Python,style=example]
import wfdbx
rec = wfdbx.read_record("100s", sampto=360)
wfdbx.write_record("demo", rec.signals, fs=rec.fs, fmt=212, gains=rec.gains,
                   baselines=rec.baselines, adczeros=1024, adcres=11,
                   units="mV", names=rec.names,
                   metadata={"Age": "69", "Sex": "M"})
print(wfdbx.read_metadata("demo"))   # [('Age', '69'), ('Sex', 'M')]
wfdbx.set_metadata("demo", "age", "70")
wfdbx.update_metadata("demo", {"Sex": None, "Dx": "Normal sinus rhythm"})
print(wfdbx.get_metadata("demo", "AGE"))         # 70
print(wfdbx.read_metadata("demo"))   # [('age', '70'), ('Dx', 'Normal sinus rhythm')]
print(wfdbx.remove_metadata("demo", "dx"))       # 1
\end{lstlisting}
After the call of \code{update\_metadata}, the header \file{demo.hea} is:
\begin{lstlisting}[style=text]
demo 2 360 360
demo.dat 212 200/mV 11 1024 995 20844 0 MLII
demo.dat 212 200/mV 11 1024 1011 28153 0 V5
# age: 70
# Dx: Normal sinus rhythm
\end{lstlisting}

\section{The functions in each language}
\label{sec:metadata-functions}

\Cref{tab:metadata-functions} lists the five functions, with their names in
each language. They take the record name and, as options, the input
directory and the network permission for reading the header, and, for the
functions that change it, the output directory \code{dir}. The chapters of
\cref{part:languages} give their exact signatures, results and errors
(\cref{chap:python,chap:cpp,chap:rust,chap:julia,chap:r,chap:matlab,chap:kotlin}).

\begin{table}[htbp]
\caption{The header metadata functions.}
\label{tab:metadata-functions}
\begin{tabularx}{\linewidth}{@{}LP{3.2cm}P{3.0cm}P{2.6cm}@{}}
\toprule
{\tabhead What it does} & {\tabhead Python, C++, Rust, Julia, R} &
  {\tabhead MATLAB/Octave} & {\tabhead Kotlin} \\
\midrule
Read all the entries, in file order, as (key, value) pairs &
  \code{read\_metadata} & \code{wfdbx.read\_metadata} &
  \code{readMetadata} \\
Return the value of one entry, or none & \code{get\_metadata} &
  \code{wfdbx.get\_metadata} & \code{getMetadata} \\
Change the value of an entry, or add it & \code{set\_metadata} &
  \code{wfdbx.set\_metadata} & \code{setMetadata} \\
Remove the entries for a key; return the number of lines removed &
  \code{remove\_metadata} & \code{wfdbx.remove\_metadata} &
  \code{removeMetadata} \\
Make several changes in one write & \code{update\_metadata} &
  \code{wfdbx.update\_metadata} & \code{updateMetadata} \\
\bottomrule
\end{tabularx}
\end{table}

The functions are called with the prefix of each language: \code{wfdbx.}
in Python and MATLAB/Octave, \code{wfdbx::} in C++ and Rust, \code{WFDBX.}
in Julia, and none in R and Kotlin. In Rust, each function also has a form
with options, \code{read\_metadata\_with}, \code{get\_metadata\_with}, and so
on. The errors are those of the high-level functions of each language
(\cref{sec:errors}). An argument error and a file error are told apart: in
Python, \code{ValueError} (\code{TypeError} for an argument of the wrong
type) and \code{OSError} (\code{FileNotFoundError} for a missing
\code{dir}); in C++, \code{wfdbx::ArgumentError} and
\code{wfdbx::FileError}; in Julia, \code{ArgumentError} and
\code{WfdbError}; in R, the condition classes
\code{wfdbx\_argument\_error} and \code{wfdbx\_file\_error}; in Kotlin,
\code{WfdbArgumentException} and \code{WfdbException}.

\begin{phibox}
Metadata often holds information about patients, such as ages, dates,
diagnoses or names. \wfdbx{} stores it as plain text, exactly as given, and
does not de-identify it, in the header or anywhere else. Removing identifying
data from headers (their metadata, and the other info strings) before records
are shared is the responsibility of the user.
\end{phibox}
